% Generated draft: requires mathematical and source review.
\documentclass[11pt]{article}
\usepackage[margin=25mm]{geometry}
\usepackage{fontspec}
\setmainfont{Latin Modern Roman}
\usepackage{amsmath,amssymb,mathtools,unicode-math}
\setmathfont{Latin Modern Math}
\usepackage{enumitem,needspace,xurl,hyperref}
\setlength{\parindent}{0pt}
\setlength{\parskip}{5pt}
\setlength{\emergencystretch}{3em}
\newcommand{\N}{\mathbb{N}}
\newcommand{\Z}{\mathbb{Z}}
\newcommand{\Q}{\mathbb{Q}}
\newcommand{\R}{\mathbb{R}}
\newcommand{\C}{\mathbb{C}}
\newcommand{\F}{\mathbb{F}}
\newcommand{\A}{\mathbb{A}}
\newcommand{\PP}{\mathbb{P}}
\newcommand{\E}{\mathbb{E}}
\newcommand{\eps}{\varepsilon}
\newcommand{\dd}{\,\mathrm{d}}
\begin{document}
\section*{UnsolvedMath ID: 30003467}
\textbf{Source problem code:} catalogue entry
\subsection{Definitions and mathematical statement}
For an explicit conventional interpretation, work in the Euclidean plane \(\R^2\). A closed Jordan region is the union of a simple closed curve and its bounded complementary component. A pseudo-disk family \(\mathcal D\) is a family of such regions satisfying
\[
\lvert\partial D\cap\partial D'\rvert\leq 2
\qquad\text{for all distinct }D,D'\in\mathcal D.
\]
This convention is provisional: the supplied record does not define pseudo-disks or specify boundary and degeneracy conventions. In this model, boundary points belong to the regions, and distinct boundaries cannot share an arc.

For a finite point set \(P\subseteq\R^2\), a three-coloring is a function \(c:P\to\{1,2,3\}\); using all three colors is not required. A set \(P\cap D\) is nonmonochromatic if it contains points receiving different colors.

The conjecture in the supplied original statement is that, for every pseudo-disk family \(\mathcal D\) and every finite point set \(P\subseteq\R^2\), there exists a three-coloring \(c\), allowed to depend on both \(P\) and \(\mathcal D\), such that
\[
\forall D\in\mathcal D,\qquad
\lvert P\cap D\rvert\geq 4
\ \Longrightarrow\
\exists p,q\in P\cap D\quad c(p)\neq c(q).
\]
Thus the threshold counts points of \(P\) in an individual member of the family. Members containing at most three points impose no condition.

The source does not explicitly restrict the family to be finite. For finite \(P\), the finite-family and arbitrary-family versions of the displayed coloring assertion are equivalent, because only finitely many distinct subsets \(P\cap D\) occur.

The supplied rewritten statement has a quantifier ambiguity: it could require one coloring of \(P\) to work simultaneously for every pseudo-disk family. The displayed formulation follows the original statement, which fixes the arrangement before choosing the coloring. The precise geometric conventions still require source verification. The supplied registry labels the problem open; that label is not an independently verified mathematical or literature assessment.
\subsection{Short English statement}
Given a pseudo-disk family in the plane and finitely many points, can the points be colored with at most three colors so that every region containing at least four of the points contains two different colors? The coloring may depend on the family. Boundary conventions and the rewritten statement's quantifier order require review.
\subsection{Sources}
\begin{itemize}
\item Local UnsolvedMath registry: supplied complete source record, permanent ID \(30003467\), title \emph{Proper Three-Colorings for Pseudo-Disk Arrangements}, exact source problem code \texttt{catalogue entry}, in the Oberwolfach Reports Open Problems collection. The supplied proposed year is 2017; no individual proposer is named.
\item Supplied report citation: \emph{Discrete Geometry (2018)}, Oberwolfach Reports, DOI 10.4171/owr/2017/19, pp.\ 1135--1205. \url{https://doi.org/10.4171/owr/2017/19}. The record describes Oberwolfach Reports as documenting mathematical workshops and their open-problem discussions. The citation's displayed year 2018 and its DOI component 2017 are retained as supplied.
\item Supplied literature source, labeled ``Keszegh and P\'alv\"olgyi on proper coloring'': \url{https://arxiv.org/abs/1612.02158}.
\item Supplied literature source, labeled ``original OWR report'': \url{https://doi.org/10.4171/owr/2017/19}.
\item Supplied literature assessment dated exactly \texttt{2026-08-21}: the record labels the problem open, describes the stronger pseudo-disk three-coloring conjecture as unresolved, mentions partial results for convex polygons, and says that a universal pseudo-disk threshold remains open. This is supplied dated curation derived from the OWR material, not an independently verified assessment. The record itself requests independent verification of time-sensitive status and cited claims.
\item The supplied statement status is \texttt{corrected\_verified}; its accompanying explanation says that mathematical content was preserved while notation, typography, and sentence structure were repaired. This is a supplied editorial qualification, not independent verification of this draft or of the cited sources.
\end{itemize}
\end{document}
